\scanpage{96}

% letterhead
INSTITUT\\
DES HAUTES ÉTUDES SCIENTIFIQUES\\
BURES-SUR-YVETTE (S.-\&-O.)\\
928-48-53

et induisent l'identité sur le gradué associé$\ldots$

Supposons maintenant $S=\operatorname{Spec} K$, $K$ un corps, et considérons un hom $i:K\longrightarrow\mathbb C$. Alors la donnée de $i$ définit le foncteur-fibre \uline{de Betti}, \deleted{soit} $F''$, à valeurs dans $\mathbb Q$.

De plus, on a des isom canoniques

\[
F''_{\mathbb C}\overset{\text{Gaga}}{\simeq} F_{\mathbb C}
\overset{\substack{\text{th. de Hodge}\\\text{des intégrales}\\\text{harmoniques}}}{\simeq}
F_{\mathbb C}^{\unsure{'|''}}
\]

D'ailleurs, l'isom $F_{\mathbb C}^{\deleted{\unsure{''}}}\simeq F'_{\mathbb C}\simeq\operatorname{gr}(F_{\mathbb C})$ est déduit de l'isom $F_{\mathbb C}\simeq F''_{\mathbb C}$, \unsure{te donnant|redonnant|traduisant} une structure \uline{réelle} sur $F_{\mathbb C}$, donnée par une involution $\sigma$, et la \uline{filtration} \inserted{can} de $F_{\mathbb C}$, par $F^{p,q}_{\mathbb C}=F^{(p)}_{\mathbb C}\cdot\sigma(F^{(q)}_{\mathbb C})$. ---
% archive sheet number: 95
